\documentclass[a4paper]{article}
\usepackage[T1]{fontenc}
\usepackage{lmodern}   % Computer Modern has no bold tt, so \bfseries is dropped
\usepackage{upquote}   % straight quotes in 8'hFF and friends
\usepackage[margin=25mm]{geometry}
\usepackage{bluespec}

\title{Typesetting Bluespec SystemVerilog with \texttt{bluespec.sty}}
\author{}
\date{}

\begin{document}
\maketitle

% How the listings are set is the document's business; style=bsv contributes
% the language and the colors, and is asked for per listing below
\lstset{basicstyle=\ttfamily\small, columns=fullflexible, keepspaces=true}

\section*{A whole file}

\lstinputlisting[style=bsv, caption={\texttt{Gcd.bsv}},
                 numbers=left, frame=single]{Gcd.bsv}

\section*{An inline listing}

A rule fires when its guard holds:
%
\begin{lstlisting}[style=bsv]
rule count (cnt < 8'd200 && !full);
   cnt <= cnt + 1;
   $display("cnt = %0h at %0t", cnt, $time);
endrule
\end{lstlisting}

\section*{Short fragments}

The literal \lstinline[style=bsv]|8'hFF| is a sized literal,
\lstinline[style=bsv]|mkReg| is a module constructor, and
\lstinline[style=bsv]|(* synthesize *)| is an attribute.

\section*{The language on its own}

Ask for \texttt{language=BSV} instead of the style and the package paints
nothing: the keyword classes, the comments, the strings and the attribute
delimiter all take whatever the document already specified. A single
\texttt{keywordstyle} reaches all five classes, because listings falls back to
it for the classes that have no style of their own.
%
\begin{lstlisting}[language=BSV, keywordstyle=\bfseries, commentstyle=\itshape]
(* synthesize *)
module mkCount (Empty);          // a monochrome document keeps its own scheme
   Reg#(Bit#(8)) cnt <- mkReg(0);
   rule count; cnt <= cnt + 1; endrule
endmodule
\end{lstlisting}

\end{document}
